\chapter{The Lebesgue Integral}
\label{ch:iii03}

The Lebesgue integral is constructed by integrating simple functions and taking monotone suprema. It is insensitive to null-set changes, linear on integrable functions, and controlled by the integral of the absolute value.

\section*{Learning goals}
After this chapter, the reader should be able to:
\begin{itemize}
\item compute Lebesgue integrals first for simple functions and then for nonnegative or integrable functions;
\item use monotone approximation by simple functions to justify integral identities;
\item separate positive and negative parts and determine when a signed function is integrable;
\item apply basic comparison and linearity properties of the integral with all finiteness assumptions checked;
\item compare Lebesgue and Riemann integrals on functions covered by both theories;
\item construct examples of measurable functions with finite, infinite, or undefined signed integrals;
\end{itemize}
\section*{Conceptual roadmap}
\[
\boxed{\text{structure}\;\longrightarrow\;\text{estimate}\;\longrightarrow\;\text{limit or representation}\;\longrightarrow\;\text{application}.}
\]

\paragraph{Mathematical setting.}
All integrands and restriction sets are measurable. Nonnegative integration uses $0\cdot\infty=0$. For a signed measurable $f$, put $f^+=\max(f,0)$ and $f^-=\max(-f,0)$. The extended integral $\int f=\int f^+-\int f^-$ is defined when at least one of these integrals is finite; integrability requires both to be finite. Arbitrary a.e. modifications require completeness to ensure measurability; equality of integrals for two already measurable representatives needs no completeness assumption.

\section{Simple-function integral}
For a nonnegative simple function $\phi=\sum a_k1_{E_k}$ on disjoint measurable level sets, define $\int\phi=\sum a_k\mu(E_k)$.


% BEGIN VOL03-EXPANSION III03-example-01
\begin{example}[A simple-function integral]\label{ex:iii03-expand-01}
On \([0,3]\), let
\[
\phi=2\,1_{[0,1]}+5\,1_{(1,3]}.
\]
The level sets are disjoint, so
\[
\int\phi=2\cdot1+5\cdot2=12.
\]
This is the defining finite-sum calculation: value times measure, summed over
measurable levels.
\end{example}
% END VOL03-EXPANSION III03-example-01
\section{Nonnegative functions}
For measurable $f\ge0$, define the integral as the supremum of simple-function integrals over $0\le\phi\le f$.

\section{Positive and negative parts}
Write $f=f^+-f^-$ with $f^+,f^-\ge0$ and $|f|=f^++f^-$.


% BEGIN VOL03-EXPANSION III03-example-02
\begin{example}[Positive and negative parts]\label{ex:iii03-expand-02}
Write \(f=f^+-f^-\) with \(f^\pm\ge0\) and \(|f|=f^++f^-\). For
\(f(x)=1/x\) for \(0<|x|<1\), with \(f(0)=0\), both \(\int f^+\) and \(\int f^-\) are infinite.
Therefore the signed Lebesgue integral is not defined by symmetric cancellation:
the expression would be the forbidden \(\infty-\infty\).
\end{example}
% END VOL03-EXPANSION III03-example-02
\section{Integrable functions}
A measurable real-valued function is integrable when $\int|f|<\infty$, which rules out the indeterminate expression $\infty-\infty$.

\section{Monotonicity}
If $0\le f\le g$, then $\int f\le\int g$ because every simple lower approximation to $f$ is also one for $g$.

\section{Linearity}
Integrable functions form a vector space and the integral is linear.

\section{Null-set invariance}
If measurable $f=g$ almost everywhere and one is integrable, then both are integrable and have the same integral.


% BEGIN VOL03-EXPANSION III03-example-03
\begin{example}[Null-set invariance]\label{ex:iii03-expand-03}
Let \(f=0\) on \([0,1]\), and let \(g\) equal \(1000\) at \(x=1/2\) and
zero elsewhere. Then \(f=g\) almost everywhere and both integrals are zero.
Lebesgue integration naturally ignores modifications on null sets, which is why
\(L^p\) spaces use almost-everywhere equivalence classes.
\end{example}
% END VOL03-EXPANSION III03-example-03
\section{Integrals over sets}
Define $\int_Ef=\int f1_E$, making restriction to measurable subsets part of the same formalism.

\section{Core structural results}

\begin{theorem}[Absolute integral inequality]\label{thm:iii03-01}
For integrable $f$, $|\int f|\le\int|f|$.
\begin{proof}
Use $-|f|\le f\le|f|$ and monotonicity of the integral.
\end{proof}
\end{theorem}

\begin{theorem}[Zero integral criterion]\label{thm:iii03-02}
If $f\ge0$ and $\int f=0$, then $f=0$ almost everywhere.
\begin{proof}
For $E_n=\{f\ge1/n\}$, the inequality $\int f\ge(1/n)\mu(E_n)$ forces every $E_n$ null; their union is $\{f>0\}$.
\end{proof}
\end{theorem}

\begin{theorem}[Linearity]\label{thm:iii03-03}
If $f,g$ are integrable and $a,b\in\mathbb R$, then $\int(af+bg)=a\int f+b\int g$.
\begin{proof}
Prove first for nonnegative simple functions on a common refinement, pass to nonnegative measurable functions by monotone approximation, then use positive/negative parts.
\end{proof}
\end{theorem}

\section{Worked examples}

\begin{example}[Integrating a simple function]\label{ex:iii03-worked-01}
Let \(s=2\mathbf1_A+5\mathbf1_B\), where \(A,B\) are disjoint measurable
sets with \(\mu(A)=3\), \(\mu(B)=1\). Then
\[
\int s\,d\mu=2\mu(A)+5\mu(B)=11.
\]
The definition is linear in the level values and uses measure only on the
corresponding level sets.
\end{example}

\begin{example}[Positive and negative parts determine integrability]\label{ex:iii03-worked-02}
For \(f=3\mathbf1_A-2\mathbf1_B\) with disjoint finite-measure sets,
\[
f^+=3\mathbf1_A,\qquad f^-=2\mathbf1_B.
\]
Thus \(\int |f|<\infty\) exactly when both \(\int f^+\) and \(\int f^-\) are
finite, and then \(\int f=\int f^+-\int f^-\).
\end{example}

\begin{example}[Indicators turn measure into integration]\label{ex:iii03-worked-03}
For measurable \(E\),
\[
\int_X\mathbf1_E\,d\mu=\mu(E).
\]
Consequently set-theoretic continuity theorems become special cases of
convergence theorems for integrals of indicators.
\end{example}

\begin{example}[A nonnegative integral can be infinite]\label{ex:iii03-worked-04}
On \((0,1)\) with Lebesgue measure, \(f(x)=1/x\) is nonnegative measurable.
For \(n\ge2\),
\[
\int_{1/n}^1\frac{dx}{x}=\log n,
\]
so monotonicity forces \(\int_0^1x^{-1}\,dx=\infty\). Lebesgue integration
allows this value for nonnegative functions without creating an undefined
\(\infty-\infty\).
\end{example}

\section{Solved dossiers}
Each dossier combines several ideas from the chapter in a longer argument. The aim is to make hypotheses, calculations, and proof strategy explicit.

\begin{problem}[Integrating a simple function]\label{prob:iii03-dossier-01}
Compute the Lebesgue integral of \(s=2\mathbf1_A+5\mathbf1_B\) for disjoint \(A,B\) with masses \(3\) and \(1\).
\end{problem}
\begin{solution}
\(\int s\,d\mu=2\cdot3+5\cdot1=11\).
\end{solution}

\begin{problem}[Positive and negative parts determine integrability]\label{prob:iii03-dossier-02}
Decompose \(f=3\mathbf1_A-2\mathbf1_B\) into positive and negative parts and compute its integral.
\end{problem}
\begin{solution}
The parts are \(3\mathbf1_A\) and \(2\mathbf1_B\); hence
\(\int f=3\mu(A)-2\mu(B)\) when both terms are finite.
\end{solution}

\begin{problem}[Indicators turn measure into integration]\label{prob:iii03-dossier-03}
Prove \(\int\mathbf1_E\,d\mu=\mu(E)\) and explain why this identifies measure as a special case of integration.
\end{problem}
\begin{solution}
It is the defining simple-function formula with one nonzero level. Thus
integrating indicators recovers the underlying measure.
\end{solution}

\begin{problem}[A nonnegative integral can be infinite]\label{prob:iii03-dossier-04}
Show \(x\mapsto1/x\) has infinite Lebesgue integral on \((0,1)\).
\end{problem}
\begin{solution}
Integrate over \([1/n,1]\): the values are \(\log n\to\infty\). By
monotonicity the full nonnegative integral is infinite.
\end{solution}

\begin{problem}[Almost-everywhere equality]\label{prob:iii03-dossier-05}
If \(f=g\) a.e. and both are integrable, prove \(\int f=\int g\).
\end{problem}
\begin{solution}
The difference vanishes a.e.; \(\int|f-g|=0\), hence
\(|\int f-\int g|\le\int|f-g|=0\).
\end{solution}

\begin{problem}[Integral monotonicity]\label{prob:iii03-dossier-06}
Prove \(0\le f\le g\) implies \(\int f\le\int g\).
\end{problem}
\begin{solution}
Approximate \(f\) from below by simple functions; each is bounded above by
\(g\), so taking suprema in the definition preserves the inequality.
\end{solution}

\begin{problem}[Absolute integrability]\label{prob:iii03-dossier-07}
Show \(|\int f|\le\int|f|\) for integrable \(f\).
\end{problem}
\begin{solution}
Since \(-|f|\le f\le|f|\), monotonicity gives
\(-\int|f|\le\int f\le\int|f|\).
\end{solution}

\begin{problem}[Finite-valued simple approximation]\label{prob:iii03-dossier-08}
Why is the integral of a nonnegative measurable function defined as a supremum over simple functions below it?
\end{problem}
\begin{solution}
Simple functions have explicit integrals, and taking the supremum over all
lower simple approximants produces the largest monotone extension compatible
with the simple-function integral.
\end{solution}

\begin{problem}[Absolute integral inequality proof dossier]\label{prob:iii03-dossier-09}
Prove the following structural statement and identify the step where its main hypothesis is used: For integrable $f$, $|\int f|\le\int|f|$.
\end{problem}
\begin{solution}
Use $-|f|\le f\le|f|$ and monotonicity of the integral. This proof also explains why the stated hypothesis is structural rather than cosmetic.
\end{solution}

\begin{problem}[Zero integral criterion proof dossier]\label{prob:iii03-dossier-10}
Prove the following structural statement and identify the step where its main hypothesis is used: If $f\ge0$ and $\int f=0$, then $f=0$ almost everywhere.
\end{problem}
\begin{solution}
For $E_n=\{f\ge1/n\}$, the inequality $\int f\ge(1/n)\mu(E_n)$ forces every $E_n$ null; their union is $\{f>0\}$. This proof also explains why the stated hypothesis is structural rather than cosmetic.
\end{solution}

\begin{problem}[Linearity proof dossier]\label{prob:iii03-dossier-11}
Prove the following structural statement and identify the step where its main hypothesis is used: If $f,g$ are integrable and $a,b\in\mathbb R$, then $\int(af+bg)=a\int f+b\int g$.
\end{problem}
\begin{solution}
Prove first for nonnegative simple functions on a common refinement, pass to nonnegative measurable functions by monotone approximation, then use positive/negative parts. This proof also explains why the stated hypothesis is structural rather than cosmetic.
\end{solution}

\begin{problem}[Chapter synthesis]\label{prob:iii03-dossier-12}
Connect the main definitions and structural theorems of this chapter into one proof strategy. Explain what object is controlled, what estimate or closure principle is used, and what conclusion becomes available.
\end{problem}
\begin{solution}
The Lebesgue integral is constructed by integrating simple functions and taking monotone suprema. It is insensitive to null-set changes, linear on integrable functions, and controlled by the integral of the absolute value. The recurring strategy is to encode the object in the chapter's natural measurable, normed, Fourier, or distributional structure, establish the controlling estimate, and then pass to the desired limit or representation.
\end{solution}
\section{Exercises with complete solutions}
The shorter exercise layer remains separate from the solved dossiers.

\begin{exercise}\label{exr:iii03-01}
State and apply the central principle behind Simple-function integral. Give one immediate consequence in the setting of this chapter.
\end{exercise}
\begin{hint}
Start with the simple-function formula \(\int \sum a_k1_{E_k}=\sum a_k\mu(E_k)\) on disjoint pieces.
\end{hint}
\begin{solution}
For a nonnegative simple function $\phi=\sum a_k1_{E_k}$ on disjoint measurable level sets, define $\int\phi=\sum a_k\mu(E_k)$. As an immediate consequence, the same principle can be inserted into later convergence, approximation, transform, or weak-form arguments whenever its hypotheses are verified.
\end{solution}

\begin{exercise}\label{exr:iii03-02}
State and apply the central principle behind Nonnegative functions. Give one immediate consequence in the setting of this chapter.
\end{exercise}
\begin{hint}
For a nonnegative function, choose an increasing sequence of simple functions and use the defining supremum/monotone limit.
\end{hint}
\begin{solution}
For measurable $f\ge0$, define the integral as the supremum of simple-function integrals over $0\le\phi\le f$. As an immediate consequence, the same principle can be inserted into later convergence, approximation, transform, or weak-form arguments whenever its hypotheses are verified.
\end{solution}

\begin{exercise}\label{exr:iii03-03}
State and apply the central principle behind Positive and negative parts. Give one immediate consequence in the setting of this chapter.
\end{exercise}
\begin{hint}
Split \(f\) into \(f^+-f^-\) and check that both positive-part integrals are finite before calling \(f\) integrable.
\end{hint}
\begin{solution}
Write $f=f^+-f^-$ with $f^+,f^-\ge0$ and $|f|=f^++f^-$. As an immediate consequence, the same principle can be inserted into later convergence, approximation, transform, or weak-form arguments whenever its hypotheses are verified.
\end{solution}

\begin{exercise}\label{exr:iii03-04}
State and apply the central principle behind Integrable functions. Give one immediate consequence in the setting of this chapter.
\end{exercise}
\begin{hint}
Use monotonicity of the integral directly when one function is bounded above by another almost everywhere.
\end{hint}
\begin{solution}
A measurable real-valued function is integrable when $\int|f|<\infty$, which rules out the indeterminate expression $\infty-\infty$. As an immediate consequence, the same principle can be inserted into later convergence, approximation, transform, or weak-form arguments whenever its hypotheses are verified.
\end{solution}

\begin{exercise}\label{exr:iii03-05}
State and apply the central principle behind Monotonicity. Give one immediate consequence in the setting of this chapter.
\end{exercise}
\begin{hint}
Linearity for signed functions is safest after verifying absolute integrability of all terms involved.
\end{hint}
\begin{solution}
If $0\le f\le g$, then $\int f\le\int g$ because every simple lower approximation to $f$ is also one for $g$. As an immediate consequence, the same principle can be inserted into later convergence, approximation, transform, or weak-form arguments whenever its hypotheses are verified.
\end{solution}

\begin{exercise}\label{exr:iii03-06}
State and apply the central principle behind Linearity. Give one immediate consequence in the setting of this chapter.
\end{exercise}
\begin{hint}
To compare with the Riemann integral, use the common simple/step approximations rather than treating the two definitions as identical.
\end{hint}
\begin{solution}
Integrable functions form a vector space and the integral is linear. As an immediate consequence, the same principle can be inserted into later convergence, approximation, transform, or weak-form arguments whenever its hypotheses are verified.
\end{solution}

\begin{exercise}\label{exr:iii03-07}
State and apply the central principle behind Null-set invariance. Give one immediate consequence in the setting of this chapter.
\end{exercise}
\begin{hint}
An infinite integral is allowed for nonnegative functions; distinguish this from an undefined expression \(\infty-\infty\).
\end{hint}
\begin{solution}
If measurable $f=g$ almost everywhere and one is integrable, then both are integrable and have the same integral. As an immediate consequence, the same principle can be inserted into later convergence, approximation, transform, or weak-form arguments whenever its hypotheses are verified.
\end{solution}

\begin{exercise}\label{exr:iii03-08}
State and apply the central principle behind Integrals over sets. Give one immediate consequence in the setting of this chapter.
\end{exercise}
\begin{hint}
For truncation, replace \(f\) by \(\min(f,n)\) or by bounded-support cutoffs and then pass to the limit.
\end{hint}
\begin{solution}
Define $\int_Ef=\int f1_E$, making restriction to measurable subsets part of the same formalism. As an immediate consequence, the same principle can be inserted into later convergence, approximation, transform, or weak-form arguments whenever its hypotheses are verified.
\end{solution}


% BEGIN VOL03-EXPANSION III03-exercises-01
\section{Graded supplementary exercises}
These exercises extend the preserved foundational set and are graded by mathematical role.

\subsection*{Standard computations}

\begin{exercise}[Simple integral]\label{exr:iii03-expand-01}
Integrate a simple function equal to 3 on a set of measure 4 and 2 on a disjoint set of measure 5.
\end{exercise}
\begin{hint}
Multiply values by measures.
\end{hint}
\begin{solution}
The integral is 22.
\end{solution}

\begin{exercise}[Indicator]\label{exr:iii03-expand-02}
What is the integral of the indicator of E?
\end{exercise}
\begin{hint}
Use the simple-function definition.
\end{hint}
\begin{solution}
It equals mu(E).
\end{solution}

\begin{exercise}[Positive part]\label{exr:iii03-expand-03}
If f is -2 on measurable A, 3 on disjoint measurable B, and zero elsewhere, identify positive and negative parts.
\end{exercise}
\begin{hint}
Keep positive magnitudes separately.
\end{hint}
\begin{solution}
f+ is 3 on B and f- is 2 on A.
\end{solution}

\begin{exercise}[Absolute inequality]\label{exr:iii03-expand-04}
State the basic absolute integral inequality.
\end{exercise}
\begin{hint}
Compare f with |f|.
\end{hint}
\begin{solution}
Absolute value of integral f is at most integral |f|.
\end{solution}

\begin{exercise}[Set integral]\label{exr:iii03-expand-05}
How is integral over E defined?
\end{exercise}
\begin{hint}
Use an indicator.
\end{hint}
\begin{solution}
Integrate f times the indicator of E.
\end{solution}

\subsection*{Proofs}

\begin{exercise}[Zero criterion]\label{exr:iii03-expand-06}
Prove nonnegative f with zero integral is zero almost everywhere.
\end{exercise}
\begin{hint}
Use level sets where f>=1/n.
\end{hint}
\begin{solution}
Each such set has measure zero; their union is where f>0.
\end{solution}

\begin{exercise}[Monotonicity]\label{exr:iii03-expand-07}
Prove 0<=f<=g implies integral f <= integral g.
\end{exercise}
\begin{hint}
Compare simple lower approximations.
\end{hint}
\begin{solution}
Every simple lower approximation to f is also one for g.
\end{solution}

\begin{exercise}[Absolute inequality]\label{exr:iii03-expand-08}
Prove the absolute integral inequality.
\end{exercise}
\begin{hint}
Use -|f| <= f <= |f|.
\end{hint}
\begin{solution}
Integrate the two inequalities and combine.
\end{solution}

\begin{exercise}[Null invariance]\label{exr:iii03-expand-09}
Prove almost-everywhere equal integrable functions have equal integrals.
\end{exercise}
\begin{hint}
Use that the integral of a measurable function supported on a null set is zero.
\end{hint}
\begin{solution}
Its integral is zero, so the difference of integrals is zero.
\end{solution}

\subsection*{Counterexamples and hypothesis testing}

\begin{exercise}[Undefined signed integral]\label{exr:iii03-expand-10}
Give a measurable signed function whose positive and negative integrals are both infinite.
\end{exercise}
\begin{hint}
Use 1/x around zero.
\end{hint}
\begin{solution}
Take $f(x)=1/x$ on $(-1,1)\setminus\{0\}$ and $f(0)=0$. The integrals of $f^+$ and $f^-$ both equal $\int_0^1 dx/x=\infty$.
\end{solution}

\begin{exercise}[Finite values not enough]\label{exr:iii03-expand-11}
Can a pointwise finite function fail to be integrable?
\end{exercise}
\begin{hint}
Use a slowly decaying tail.
\end{hint}
\begin{solution}
Yes; 1/x on [1,infinity) is finite pointwise but has infinite integral.
\end{solution}

\begin{exercise}[Lebesgue not Riemann]\label{exr:iii03-expand-12}
Give a bounded Lebesgue-integrable function on [0,1] that is not Riemann integrable.
\end{exercise}
\begin{hint}
Use the rational indicator.
\end{hint}
\begin{solution}
Take $f=1_{\mathbb Q\cap[0,1]}$. Its Lebesgue integral is zero, but every nondegenerate subinterval has infimum zero and supremum one; lower and upper Riemann integrals are zero and one.
\end{solution}

\subsection*{Applications and investigations}

\begin{exercise}[Expectation]\label{exr:iii03-expand-13}
Why is expectation a Lebesgue integral?
\end{exercise}
\begin{hint}
Probability is a measure and a random variable is measurable.
\end{hint}
\begin{solution}
Expectation is integral of the random variable against probability.
\end{solution}

\begin{exercise}[Sparse corruption]\label{exr:iii03-expand-14}
For Lebesgue measure on the real line, why can changing finitely many sample values leave an integral unchanged?
\end{exercise}
\begin{hint}
A finite set has measure zero.
\end{hint}
\begin{solution}
Null-set changes do not affect the integral.
\end{solution}

\subsection*{Challenge problems}

\begin{exercise}[Layer-cake]\label{exr:iii03-expand-15}
On a sigma-finite measure space, explain the level-set representation of a nonnegative measurable integral.
\end{exercise}
\begin{hint}
Represent f(x) as the integral over t of the indicator f(x)>t.
\end{hint}
\begin{solution}
Tonelli yields integral f = integral over t of the measure of {f>t}.
\end{solution}

\begin{exercise}[Power singularity]\label{exr:iii03-expand-16}
When is x\textasciicircum{}(-alpha) integrable on (0,1)?
\end{exercise}
\begin{hint}
Integrate the power near zero.
\end{hint}
\begin{solution}
Exactly when alpha<1.
\end{solution}
% END VOL03-EXPANSION III03-exercises-01
\section*{Chapter summary}
The chapter now supplies a canonical theorem-and-dossier layer for subsequent measure, Fourier, distributional, Sobolev, and PDE arguments.
